\reviewsection{Meaningful Communication Is Authorship, Not Performance}

Meaningful communication is not defined by speech, speed, eye contact, independent device use, or whether a response matches an adult's expectation. It exists when a student's message is recognized as theirs, understood or respectfully repaired, and able to influence learning, relationships, support, and what happens next. Augmentative and alternative communication (AAC) can make that authorship more accessible, but the presence of a device does not prove that the student has communication power. AAC may include gestures, facial expressions, eye gaze, writing, drawing, symbols, partner-assisted scanning, low-tech boards, applications, and speech-generating devices. There is no age, cognitive score, or developmental milestone that a student must reach before communication support can help \citep{ashaAAC}.

\citet{norrie2021context} make the difference between device presence and communication access visible. In their five-month study of one special-education school, only six of approximately 180 pupils with complex communication needs were active high-tech AAC users. Their findings locate barriers not only in a device or learner but in training, technical support, staffing, policy, infrastructure, and the fit between technology and daily life. The study changes the question from ``Why is the student not using the device?'' to ``What conditions make this communication system usable across real activities and partners?''

That distinction matters because a learner can technically possess AAC while having access to only a narrow adult agenda. A system limited to requesting food, answering teacher questions, or following directions does not provide the language needed to disagree, joke, ask an unexpected question, express identity, repair a misunderstanding, participate academically, or refuse. Meaningful communication must include functions that adults did not script in advance.
